\cotefile{149}

\scanpage{5}
\mspage{2}
\margin{\diagram{\begin{tikzcd}[ampersand replacement=\&,column sep=small,row sep=small]
K\arrow[r,dash]\arrow[dr,dash]\&K'\arrow[dr,dash]\&{}\\
{}\&\bar K\arrow[r,dash]\&\bar K'
\end{tikzcd}}}
$K_1$ de $\bar K/K$, clôture \inserted{alg} sép. de $K$ dans $K'$, i.e. $K_1=\bar K\cap K'$.

Quand $K'$ est une extension de type fini de $K$, $K_1$ est une ext. finie de $K$, et on en conclut que l'image de (3) est bien un sous-groupe d'indice fini, égal : $E_{\bar K/K}$ sss \inserted{$K_1=K$ i.e.} $K$ est séparablement alg. clos dans $K'$. D'ailleurs, on montre sans mal que \inserted{(si $K'$ est ext. de type fini)} l'hom (3) est injectif sss $K'$ est une extension algébrique de $K$. Donc il est bijectif sss $K'$ est une extension radicielle de $K$. Dans la suite nous nous bornons (provisoirement) aux corps de car. $0$, et la condition précédente signifie \inserted{donc} que $K\unsure{-|\longrightarrow}K'$ est un \uline{isomorphisme}.

Ainsi, le foncteur $K\longrightarrow B_K$ ou $K\unsure{-|\longrightarrow}\Pi_K$, ou $(K,\bar K)\unsure{-|\longrightarrow}E_{\bar K/K}$, est \uline{conservatif} quand on se limite comme morphismes de corps \inserted{$K\longrightarrow K'$ (de car $0$)} à ceux \deleted{\unsure{pour}} \unsure{qui|que} font de $K'$ une ext. de type fini de $K$\unsure{.|,}
\margin{4}

\scanpage{6}
P.ex. il suffit de se limiter aux extensions de type fini du corps premier $\mathbb{Q}$ -- on trouve un foncteur conservatif de la cat. de ces corps dans celle des groupoïdes \inserted{(ou des topos)}, au sens : un morphisme de corps qui donne une équivalence de groupoïdes (ou de \unsure{topos|Topos}) est un \uline{iso}.

\margin{en car $p>0$ le foncteur est conservatif si de plus on passe à la catégorie quot$^t$ de celle des corps\unsure{,} où une extension radicielle finie est considérée comme triviale... (le foncteur $K-B_K$ ou $K-\Pi_K$ se factorise par cette cat. quotient\unsure{..|...})}

Quand on prend des corps quelconques, le 2-foncteur $K-B_K$ ou $K\longrightarrow\Pi_K$ \inserted{ou $(K,\bar K)\unsure{-|\longrightarrow}E_{\bar K/K}$} est cependant loin d'être fidèle. Ainsi, si $K$ est sép clos, $B_K$ est le ``topos ponctuel'', $\Pi_K$ le groupoïde ponctuel, $E_{\bar K/K}\simeq(1)$ -- il est clair que les morphismes entre corps sép. clos ne sont pas décrits par les morphismes entre leurs topos étales, ou groupoïdes fondamentaux ! Pour cette raison, il y a lieu d'associer à un corps $K$ un objet plus fin que $B_K$ ou $\Pi_K$,
\margin{5}

\scanpage{7}
\mspage{3}
à savoir le système projectif des $B_{K_i}$, ou des $\Pi_{K_i}$, pour $K_i$ sous-corps de $K$ de type fini sur le corps premier, et à un syst. $(K,\bar K)$ le système projectif des $E_{\bar K_i/K_i}$, où $\bar K_i$ est la clôture alg. sép. de $K_i$ dans $\bar K$. On aura
\[
(4)\qquad\left\{\begin{aligned}
\Pi_K&\cong\varprojlim\Pi_{K_i}\\
B_K&\cong\varprojlim B_{K_i}\\
E_{\bar K/K}&\xrightarrow{\sim}\varprojlim E_{\bar K_i/K_{\unsure{i}}}
\end{aligned}\right.
\]
i.e. on réconstitue les objets $B_K$, $\Pi_K$, $E_{\bar K/K}$ à partir du syst. proj. correspondant -- mais l'inverse n'est pas vrai.

En fait, comme le foncteur $\operatorname{Ind}(\text{corps.t.f})\longrightarrow\text{corps}$ de la catégorie des syst. inductifs de corps de type fini, vers celle des corps, est une équivalence de catégories (pour des raisons triviales), il s'ensuit que le foncteur $K\longrightarrow\mathbb{B}_K$, ou $K\longrightarrow\boldsymbol{\Pi}_K$, ou $(K,\bar K)\longrightarrow\mathbb{E}_{\bar K/K}$, \unsure{allant|elles}
\margin{6}

\scanpage{8}
des corps vers les proobjet idoines, \unsure{va|vont} avoir toutes les propriétés de fidélité du foncteur $K-B_K$, ou $K-\Pi_K$, ou $(K,\bar K)-\Gamma_{\bar K/K}$, \uline{restreint} aux corps \inserted{absolument} de type fini, auxquels nous allons par la suite nous borner, la plupart du temps. Mais il sera nécessaire, au cours du travail, de donner une description \inserted{purement} algébrique, p.ex., du pro-pro-groupe \unsure{fini} associé p.ex. à $\mathbb C$ (plus précisément, à $(\mathbb C,\mathbb C)$ !).

Le rôle dominant sera joué par le corps premier de car. $0$, $\mathbb Q$, donc par $B_{\mathbb Q}$ et $\Pi_{\mathbb Q}$, qui a un objet canonique, noté $\overline{\mathbb Q}_0$ -- la clôture alg. de $\mathbb Q$ dans $\mathbb C$. On pose.
\[
(5)\qquad \Gamma_{\mathbb Q}=E_{\overline{\mathbb Q}_0/\mathbb Q}
\]
\margin{et on écrira \inserted{souvent} $\Gamma_{\overline{\mathbb Q}/\mathbb Q}$ au lieu de $E_{\overline{\mathbb Q}/\mathbb Q}$, pour une clôture alg. quelq. $\overline{\mathbb Q}$ de $\mathbb Q$}
Pour tt corps $K$ de car. $0$ -- en particulier pour les corps $K$ de type fini sur $\mathbb Q$, aux-
\margin{7}

\scanpage{9}
\mspage{4}
quels nous allons nous borner par la suite -- on a donc des hom canoniques
\[
(6)\qquad B_K\unsure{-|\longrightarrow}B_{\mathbb Q}\qquad\Pi_K\longrightarrow\Pi_{\mathbb Q}
\]
qui s'explicitent, quand on a choisi un objet de $\Pi_K$ i.e. un $\bar K/K$, d'où un $\overline{\mathbb Q}/\mathbb Q$, par un hom de groupes profinis
\[
(7)\qquad E_{\bar K/K}\longrightarrow\Gamma_{\overline{\mathbb Q}/\mathbb Q}.
\]
Par la suite, on regardera toujours $B_K$, $\Pi_K$ ou $E_{\bar K/K}$ comme munis de cette structure supplémentaire -- ce sont les \deleted{f}morphismes (de topos, de groupoïdes, ou de groupes profinis) ``\uline{arithmétiques}'', domineront la situation. !

Un intérêt particulier s'attache au noyau de (7), que je note $\pi_{\bar K/K}$ -- ``\unsure{on|en} l'appelle ``partie géométrique'' du groupe Galois $E_{\bar K/K}$, par opposition au quotient
\[
E_{\bar K/K}/\pi_{\bar K/K}=\Gamma^{\unsure{\deleted{a}|a}}_{\bar K/K}\hookrightarrow\Gamma_{\overline{\mathbb Q}/\mathbb Q},
\]
\margin{On va noter $\Gamma=\Gamma^{\unsure{\deleted{a}|\deleted{can}|can}}_{\bar K/K}$ cette ``partie arithmétique''}
que j'appelle sa
\margin{8}

\scanpage{10}
partie ``arithmétique'' -- celle-ci est un sous-groupe ouvert de $\Gamma_{\overline{\mathbb Q}/\mathbb Q}$, qui, on \unsure{l'a|le} vu, correspond au sous-corps $K_0$ de $\overline{\mathbb Q}/\mathbb Q$, \unsure{u|``} extension finie $/\mathbb Q$ de $\overline{\mathbb Q}/\mathbb Q$, clôture alg de $\mathbb Q$ dans $K$, de sorte qu'on a une suite exacte

\noindent (8)\quad\diagram{\begin{tikzcd}[ampersand replacement=\&,column sep=small,row sep=small]
1\arrow[r,dash]\&\pi_{\bar K/K}\arrow[r]\&E_{\bar K/K}\arrow[r]\&\unsure{\Gamma|\Gamma'}_{\bar K/K}\arrow[r]\arrow[d,hook]\&1\\
{}\&{}\&{}\&\Gamma_{\overline{\mathbb Q}/\mathbb Q}\&{}
\end{tikzcd}}

\margin{\uline{NB} $\pi_{\bar K/K}=(1)$ sss $K$ alg. sur $\mathbb Q$, i.e. fini sur $\mathbb Q$}
On va donner une interprétation de ce noyau, et de la suite exacte (8), en écrivant
\[
(9)\qquad K=\varinjlim_i A_i
\]
où les $A_i$ sont les sous-$\mathbb Q$-algèbres de type fini de $K$, correspondant au syst. projectif des ``modèles affines'' $U_i=\operatorname{Spec}(A_i)$ de $K/\mathbb Q$. Parmi les $A_i$, il y a d'ailleurs un syst. cofinal formé de $A_i$ réguliers, \inserted{i.e. les $U_i$ ``lisses$/\mathbb Q$,} avec comme morphismes de transition des morphismes de localisation $A_i\longrightarrow A_j=(A_i)_{f_{ij}}$ ($f_{ij}\in A_i\setminus\{0\}$) -- donc les $U_j\lhook\joinrel\relbar U_i$ des inclusions ouvertes, $U_j=(U_i)_{f_{ij}}$. On peut
\margin{9}

\scanpage{11}
\mspage{5}
même, d'après Mike Artin, prendre pour $U_i$ des schémas ``élémentaires'' sur $K_0$, se dévissant en fibrations successives de courbes. \inserted{Notons que $\operatorname{Spec}K=\eta$ est le p$^t$ générique} \deleted{Ceci dit, le \unsure{pt|dt|ch}} commun des $U_i$, qui sont géom$^t$ conn. \inserted{(et $\unsure{\simeq|=}$ géom. intègres)} sur \deleted{$K_0$} $k$ (clôture alg. de $\mathbb Q$ dans $K$).

Le choix de $\bar K$ définit un pt géom. $\bar\eta$ sur les $U_i$, d'où des groupes $\pi_1(U_i,\bar\eta)=E_i$, et on a, en vertu des résultats bien connus
\[
\operatorname{Spec}K=\varprojlim U_i
\]
\[
(10)\qquad E_{\bar K/K}=\pi_1(\eta,\bar\eta)\xrightarrow{\sim}\varprojlim_i\deleted{\unsure{\pi}}E_i
\]
\margin{$(\Gamma_i=\pi_1(U_i,\bar\eta))$}
D'autre part, si on pose
\[
(11)\qquad\bar U_i=U_i\otimes_k\overline{\mathbb Q}
\]
on a \inserted{pour tt $i$} une suite exacte d'homotopie

\noindent (12)\quad\diagram{\begin{tikzcd}[ampersand replacement=\&,column sep=small,row sep=small]
1\arrow[r]\&\underbrace{\pi_1(\bar U_i,\bar\eta)}_{=\pi_i}\arrow[r]\&\pi_1(U_i,\bar\eta)\arrow[r]\arrow[d,equal]\&\pi_1(k,\bar\eta)\arrow[r]\arrow[d,equal]\&1\\
{}\&{}\&E_i\&\Gamma_{\bar k/k}=\Gamma_{\overline{\mathbb Q}/k}=\Gamma'_{\deleted{\unsure{\overline{\mathbb Q}/k|\bar k/k}}}\&{}
\end{tikzcd}}

\margin{\diagram{\begin{tikzcd}[ampersand replacement=\&,column sep=small,row sep=small]
{}\&\operatorname{Spec}(K)=\eta\arrow[r,dash]\arrow[d]\&\eta_{\overline{\mathbb Q}}\arrow[r,dash]\arrow[d,dash]\&\bar\eta=\operatorname{Spec}(\bar K)\\
{}\&U_i\arrow[r,dash]\arrow[d,dash]\&U_{i\overline{\mathbb Q}}=\bar U_i\arrow[d,dash]\&{}\\
\mathbb Q\arrow[r,dash]\&k\arrow[r,dash]\&\overline{\mathbb Q}\&{}
\end{tikzcd}}}
qui forment un syst. projectif de suites exactes, ou d'extensions ayant toutes même quotient $\Gamma'$, et dont les noyaux
\[
\pi_i=\pi_1(\bar U_i,\bar\eta)
\]
\margin{10}

\scanpage{12}
sont des groupes fondamentaux ``géométriques'' --- \unsure{qui|que} peuvent d'ailleurs se calculer par voie transcendante, en utilisant un plongement de $\overline{\unsure{\deleted{\mathbb Q}K|\deleted{K}\mathbb Q}}$ dans $\mathbb C$ (d'où un isom $\overline{\mathbb Q}\simeq\overline{\mathbb Q}_0$), comme les complétés profinis des $\pi_1(U_i(\underset{\unsure{\mathbb Q|k|\mathbb C}}{\mathbb C}),\bar\eta)$, \unsure{en|où} maintenant $\bar\eta$ est interprété comme un pt \deleted{\unsure{complexe|complexes}} commun des espaces top. sous-jacents aux variétés complexes $U_i(\mathbb C)$.

La suite exacte (8) est donc la limite projective des suites exactes d'homotopie (12), ce qui donne en particulier
\[
\begin{aligned}
(13)\qquad \pi_{\bar K/K}&\simeq\varprojlim_i\underbrace{\pi_1(\bar U_i,\bar\eta)}_{\pi_i}\\
&\simeq\varprojlim_i\pi_1(U_i(\mathbb C),\bar\eta)^{\wedge}
\end{aligned}
\]
\margin{Cette interprétation n'est évidemment pas particulière au corps de base $\mathbb Q$ --- elle vaut \unsure{pour|pr} tte ext. de corps $K/K_0$}

Utilisant les fibrations des $U_i$ (dans le cas où on s'astreint à prendre des variétés élémentaires d'Artin), on trouve que tout $\pi_i$ est un groupe \deleted{\unsure{par|pr}} extension successive de \inserted{$n$} groupes profinis \uline{libres} (\unsure{ici|où} $n=\dim U_i=\operatorname{deg.tr.}K/\mathbb Q$). Ceci redonne
\margin{11}

\scanpage{13}
\mspage{6}
p.ex. que la dim coh. de $\pi_i$ est égale à $n$, celle de $E_i$ est $\leq n+2$ (pour des coeff. \deleted{$\neq2$} de $m$-torsion, quand $m$ est impair) --- et par passage à la limite, des majorations correspondantes pour les dim coh. de $\pi_{\bar K/K}$ et de $E_{\bar K/K}^{\unsure{\ell|t}}$
\[
(\unsure{14})\qquad \dim\operatorname{coh}\pi_{\bar K/K}\leq n,\qquad
\dim\operatorname{coh}_{\ell}E_{\bar K/K}\leq n+2\quad\text{si }\ell\neq2
\qquad(\ell\text{ nb premier})
\]
qui sont en fait même des \uline{égalités} (\unsure{sauf|s.f.} erreur), et \unsure{donnent|donnant} donc une description cohomologique simple du degré de transcendance absolu de $K$.

\uline{Théorème 1} Soit $K$ un corps \inserted{extension} de type fini de $\mathbb Q$. \deleted{Alors} $\bar K$ une clôture algébrique de $K$. Alors pour tout sous-groupe ouvert $\unsure{\mathcal E|\varepsilon}$ de $E_{\bar K/K}$, son centralisateur dans $E_{\bar K/K}$ est réduit au groupe unité. Itou pour $\pi_{\bar K/K}$.

Démonstration. Soit \deleted{$\unsure{E'\subset E}$}\inserted{$\Gamma'\subset\Gamma\subset$}$\subset\Gamma_{\overline{\mathbb Q}/\mathbb Q}$ l'image de $E$ dans $\Gamma$\inserted{$=\Gamma_{\bar K/K}$}, qui est donc un ss-groupe ouvert. \inserted{L'image dans $\Gamma$ du centralisateur de $\unsure{\mathcal E|\varepsilon|\mathcal E'|\varepsilon'}$ dans $\unsure{\mathcal E|\varepsilon|\mathcal E'|\varepsilon'}$ est contenue dans \unsure{le} centralisateur de $\Gamma'$ dans $\Gamma$\unsure{.|,}} \deleted{Si $g\in\Gamma_{\bar K/K}$ centralise $\unsure{\mathcal E|\varepsilon}$, son image $\gamma$ dans $\Gamma$ centralise $\Gamma'$}\unsure{\deleted{$\unsure{\subset|\text{de}|\text{et}}\Gamma_{\mathbb Q}$}|$\unsure{\subset|\text{de}|\text{et}}\Gamma_{\mathbb Q}$} Je dis qu'il est
\margin{12}

\scanpage{14}
égale à $1$, ce qui équivaut donc au

\uline{Corollaire} Dans $\Gamma_{\mathbb Q}=\Gamma_{\overline{\mathbb Q}_0/\mathbb Q}$, le centralisateur de tout sous-groupe ouvert est réduit à $(1)$.

OPS le sous-groupe ouvert \inserted{$\Gamma'$} invariant, Il est bien connu (?) que son centre est réduit à $1$
\margin{? à vérifier}
donc si $\unsure{Z|\mathcal Z}$ est son centralisateur dans $\Gamma_{\mathbb Q}$, l'hom $\unsure{Z|\mathcal Z}\longrightarrow\Gamma_{\mathbb Q}/\Gamma'$ est injectif donc $\unsure{Z|\mathcal Z}$ est fini. Mais on sait que les seuls \unsure{éléments|élémts} $\neq1$ de $\Gamma_{\mathbb Q}$ d'ordre fini sont les conjugués de $\tau$, conjugaison complexe. Mais le centralisateur de $\tau$ dans $\Gamma_{\overline{\mathbb Q}}$ est réduit à $(1,\tau)$, donc ne peut contenir $\Gamma'$, donc $\tau\notin\unsure{Z|\mathcal Z}$, donc $\unsure{Z|\mathcal Z}=(1)$.

Revenant à $\unsure{\mathcal E|\varepsilon}\subset E_{\bar K/K}$, on voit donc que son centralisateur $\unsure{Z|\mathcal Z}$ dans $E_{\bar K/K}$ a une image dans $\Gamma$ \unsure{réduite|réduit} à $\{1\}$, donc $\unsure{Z|\mathcal Z}\subset\pi_{\bar K/K}$. Soit $\pi'\subset\pi=\pi_{\bar K/K}$ la trace de $\unsure{\mathcal E'|\varepsilon'}$ sur $\pi$, c'est un sous-groupe ouvert de $\pi$, et on est ramené à
\margin{13}

\scanpage{15}
\mspage{7}
voir que $\operatorname{Centr}_{\pi}(\pi')=\{1\}$, i.e. le

\uline{Corollaire} Soit $\pi$ un groupe profini, extension successive de groupes profinis libres. Alors le centralisateur \inserted{$\unsure{Z|\mathcal Z}$ dans $\pi$} de tout sous-groupe ouvert $\pi'$ de $\pi$ est réduit à $\{1\}$.

Par dévissage on est ramené au cas d'un groupe \inserted{profini} \uline{libre}. \inserted{On sait que $\pi'$ est donc libre.} OPS $\pi'$ invariant, et on admet que \unsure{\deleted{son}le|le|les} centre d'un groupe profini libre est réduit à $1$.
\margin{à vérifier}
Donc $\unsure{Z|\mathcal Z}\longrightarrow\pi/\pi'$ est injectif, donc $\unsure{Z|\mathcal Z}$ est fini, et on admet que dans un groupe profini libre, il n'y a pas d'élément d'ordre fini $\neq1$ \margin{à vérifier} --- ce qui achève la démonstration.

\uline{Scholie} Le fait que $E_{\bar K/K}$ soit \deleted{sans}\inserted{à} \deleted{\unsure{torsion|tors}\unsure{:|.}}\inserted{centre} trivial peut s'exploiter, en disant que le groupoïde $\Pi_{\deleted{\bar K/}\unsure{K|\deleted{K}}}$ \inserted{(\unsure{ou le|où le} topos $B_K$)} est connu à équivalence près, \deleted{\unsure{que l'on|question|équivalence}}\inserted{\unsure{defini|definie} à isom} unique près, quand on connaît le groupe \uline{extérieur} associé à $E_{\bar K/K}$. Les hom
\margin{14}

\scanpage{16}
$E_{\bar K'/K'}\longrightarrow E_{\bar K/K}$ associés à des hom \inserted{$K\longrightarrow K'$} d'extensions de type fini de $\mathbb Q$, ayant une image ouverte donc un centralisateur réduit à $1$, on voit de $=$ que l'hom de topos $B_{K'}-B_K$ \unsure{ou|en|on} de groupoïdes $\Pi_{K'}-\Pi_K$, sont déterminés à \deleted{\unsure{iso}\illegible{}}\unsure{équivalence} près (définis à isom unique près) par l'hom correspondant des groupes extérieurs. Il en \unsure{sera|dc.} en particulier ainsi du morphisme \unsure{structural|structurel} $B_K\longrightarrow B_{\mathbb Q}$ \unsure{ou|en|on} $\Pi_K-\Pi_{\mathbb Q}$, qu'on peut interpréter \inserted{intrinsiquement} comme un hom de groupes \inserted{profinis} extérieurs $E_K-E_{\mathbb Q}$. Mais nous préférons suivre un point de vue un peu \unsure{affiné|différent}, en exploitant le fait que $\pi_{\bar K/K}$ est lui aussi à centre trivial. Cela signifie que l'extension de $\unsure{E|\Gamma}=\unsure{\Gamma|E}_{\bar K/K}$ par $\pi_{\bar K/K}$ est entièrement connue, \inserted{à isom unique près,} pour $\pi_{\bar K/K}$ \inserted{et $\Gamma$} fixés,
\margin{15}

\scanpage{17}
\mspage{8}
en termes de \deleted{\unsure{ain|un|une}} l'\uline{action extérieure} correspondante de $\Gamma$ sur $\pi$, comme l'image inverse de l'extension universelle
\noindent (15)\quad\diagram{\begin{tikzcd}[ampersand replacement=\&]
1\arrow[r]\&\pi\arrow[r]\&\operatorname{Aut}(\pi)\arrow[r]\&\operatorname{Aut\,ext}(\pi)\arrow[r]\&1\\
\&\&\&\Gamma\arrow[u,"\varphi"']\&
\end{tikzcd}}

Pour $K$ fixé, donc $k$ fixé, dire qu'on fixe un $\Gamma=\Gamma_{\overline{\mathbb Q}/k}$ revient à dire qu'on fixe une clôture alg. de $k$, dire qu'on fixe \inserted{un} $\pi_{\bar K/K}=\pi_1(K\otimes_k\bar k)$ signifie \inserted{est} qu'on fixe un rev. universel de $\operatorname{Spec}(K\otimes_k\bar k)=\eta\otimes_k\bar k$, \deleted{\unsure{compatible|correspond}} les deux ensembles reviennent à se donner le rev. universel \deleted{$S\bar K$} $\bar\eta=\operatorname{Spec}(\bar K)$ de $K$. Par la suite, nous décrivons (avec une fidélité qui met à cœur) les couples $(K,\bar K)$ d'une ext.\inserted{$K$} de $\mathbb Q$ de type fini, et d'une clôture algébrique $\bar K$ de $K$, par les triples $(\pi,\Gamma,\varphi)$, où $\pi=\pi_{\bar K/K}$ \inserted{et $\Gamma=\Gamma_{\bar K/K}^{\unsure{can}}$} sont des groupes profinis, et $\varphi:\Gamma\unsure{-|\longrightarrow}\operatorname{Aut\,ext}(\pi)$ une action extérieure de $\Gamma$ sur $\pi$ --- ce qui
\margin{16}

\scanpage{18}
permet de réconstituer l'extension $E_{\bar K/K}$ de $\Gamma^{\unsure{a|\infty|can}}_{\bar K/K}$ par $\pi_{\bar K/K}$. J'ai oublié de préciser qu'il faut \uline{de plus} se donner $\Gamma$ comme sous-groupe d'un $\Gamma_{\overline{\mathbb Q}/\mathbb Q}$ bien déterminé, i.e. qu'il faut se donner un objet de $\Pi_{\mathbb Q}$ et une action fidèle de $\Gamma$ dessus --- pour reconstruire par l'\unsure{ens|en} de ces données un hom de groupoïdes profinis $\Pi_K-\Pi_{\mathbb Q}$, plus un objet de $\Pi_K$ --- ou encore, un morphisme de topos pro\unsure{-|.}galoisiens $B_K\longrightarrow B_{\mathbb Q}$, plus un pt de $B_K$.

On peut aussi fixer un objet de $\Pi_{\mathbb Q}$, i.e. un pt de $B_{\mathbb Q}$, i.e. un $\overline{\mathbb Q}$, et étudier les $K$, \deleted{\unsure{Ainsi|tels} que $k$}\inserted{avec} un plongement de $k$ (clôture alg. de $\mathbb Q$ dans $K$) dans $\overline{\mathbb Q}$ --- \deleted{ils sont} mais sans se donner une clôture alg. $\bar K$ de $K$ qui induise $\overline{\mathbb Q}/\mathbb Q$. Ils sont décrits
\margin{17}

\scanpage{19}
\mspage{9}
On a ainsi plusieurs points de vue essentiellement équivalents, pour décrire par voie profinie \unsure{\deleted{\unsure{une|d'une|des}}|\unsure{une|d'une|des}} extension \inserted{$K$} de type fini de $\mathbb Q$ :

1) Par le topos étale $B_K$, en tant que topos pro-galoisien sur $B_{\mathbb Q}$

2) Par le groupoïde fondamental $\Pi_K$ de \unsure{ce|en} topos, (le groupoïde de ses pts, ou de ses rev. universels) --- en tant que groupoïde au-dessus de $\Pi_{\mathbb Q}$

\deleted{2)} \deleted{\unsure{2}}3) Par le groupe extérieur $\unsure{\mathcal E|E}_K$, \unsure{au-dessus|--- disons|--- dessus} du groupe extérieur $\unsure{\mathcal E|E}_{\mathbb Q}\longleftarrow\Gamma_{\mathbb Q}$ (pro-gr)

4) En termes d'une clôture algébrique $\bar K/K$ (i.e. en décrivant le couple $(K,\bar K)$ plutôt que $K$), par un objet $\overline{\mathbb Q}\in\operatorname{Ob}\Pi_{\mathbb Q}$\unsure{,|)} et un hom de groupes profinis $\unsure{\mathcal E|E}\unsure{\longrightarrow|-}\Gamma_{\overline{\mathbb Q}/\mathbb Q}$.
\deleted{un ss-groupe ouvert $\Gamma$ de}

5) En termes d'une clôture alg. \inserted{\unsure{fixée|fixe}} $\overline{\mathbb Q}$ de $\mathbb Q$, \deleted{par le groupe}\inserted{et \unsure{si} $\Gamma=\Gamma_{\overline{\mathbb Q}/\mathbb Q}$}, \deleted{par} considérant les couples $(K,i)$ où $i:k\longrightarrow\overline{\mathbb Q}$ est un plongement de le \unsure{clôture|cloture} alg.\inserted{$k$} de $\mathbb Q$ dans $K$ dans $\overline{\mathbb Q}$ : \deleted{\unsure{comme|connexe|connexes}} par le groupe extérieur $\pi_K\unsure{\neq|=|\simeq}\pi_1(K)$, sur lequel un sous-groupe ouvert $\Gamma_K\subset\Gamma$ opère extérieurement.

\margin{6) En termes d'un $\overline{\mathbb Q}/\mathbb Q$ : par le \unsure{gr.-oïde|gr.-oïdes} $\Pi_{K\otimes_{\mathbb Q}\overline{\mathbb Q}}$ (pas néc. connexe) sur lequel $\Gamma_{\overline{\mathbb Q}/\mathbb Q}$ opère --- ou bien \unsure{plutôt|plutot}, par le topos quotient $B_{K\otimes_{\mathbb Q}\overline{\mathbb Q}}$\unsure{,|)} sous $\Gamma_{\overline{\mathbb Q}/\mathbb Q}$ opérant dessus\ldots}
\margin{18}

\scanpage{20}
par des groupes \inserted{profinis} \uline{extérieurs} $\pi_1(K)=\Gamma_K$, sur lesquels \inserted{un sous-groupe ouvert $\Gamma$ (\unsure{non-précisé|non précisé} d'avance) de} $\Gamma_{\overline{\mathbb Q}/\mathbb Q}$ opère \uline{extérieurement}.

Un hom de corps $K\unsure{-|\longrightarrow}K'$ donne lieu à un hom de groupes extérieurs,
\margin{\unsure{on|en} suppose pour simplifier que c'est relatif \unsure{au} $=$ $k\subset\overline{\mathbb Q}$\ldots}
\deleted{$E'\longrightarrow E$} $\pi'\longrightarrow\pi$, où l'image de $\pi'$ dans $\pi$ est ouverte donc de centralisateur réduit à $(1)$, ce qui \deleted{signifie}\inserted{implique} encore que le morphisme de topos $B_{K'\otimes_k\overline{\mathbb Q}}\longrightarrow B_{K\otimes_k\overline{\mathbb Q}}$ est déterminé (à isom unique près) par cet hom extérieur. De plus on a des actions extérieures de \inserted{$\Gamma=$}$\Gamma_K$\inserted{$\subset\Gamma_{K'}$} sur $\pi'$ et $\pi$, de façon que $\pi'-\pi$ y commute --- et ceci suffit pour réconstituer, d'une part les groupes extérieurs $\mathcal E$, $\mathcal E'$ extensions \inserted{(``extérieures'')} de $\Gamma$ par $\pi$, $\pi'$ (et, à équivalence rigide près, les $B_K$, $B_{K'}$ et $B_K-B_{\mathbb Q}$, $B_{K'}-B_{\mathbb Q}$) et de plus l'hom d'extensions \inserted{extérieures} $\mathcal E\longrightarrow\mathcal E'$ de $\Gamma$.
\margin{19}

\scanpage{21}
\mspage{10}
\uline{Remarque} Quand $\pi\neq(1)$, \inserted{i.e. $K$ pas fini sur $\mathbb Q$,} le théorème 1 peut se renforcer, \inserted{sauf erreur,} en \unsure{écrivant|établissant|évidemment} que pour tt ss-groupe \inserted{$\pi'\subset\pi$} ouvert \uline{dans $\pi$}, $\operatorname{Centr}_{\unsure{\mathcal E|E|\varepsilon}}(\pi')=\{1\}$. Si $\unsure{Z|\mathcal Z}$ est ce centralisateur, on a $\unsure{Z|\mathcal Z}\cap\pi=(1)$ d'après le \unsure{th.|Th.|H.}1, \deleted{et}\inserted{prouvons que} l'image de $\unsure{Z|\mathcal Z}$ dans $\Gamma_{\bar K/K}\subset\Gamma_{\overline{\mathbb Q}/\mathbb Q}$ est \deleted{\unsure{contenue|contenu} dans le \unsure{noy}}\inserted{finie} (ce qui prouvera \inserted{\unsure{déjà}} \unsure{alors|dans|donc}, que $\unsure{Z|\mathcal Z}$ est d'ordre $1$ ou $2$, et dans le deuxième cas que son image dans $\Gamma_{\overline{\mathbb Q}/\mathbb Q}$ est \deleted{\unsure{un|ou} \unsure{él.|$(1)$, $Z$|$(1)$, $\mathcal Z$}} engendrée par un $\tau$ de conjugaison complexe) --- Quitte à remplacer $\unsure{\mathcal E|E|\varepsilon}$ par un sous-groupe ouvert assez petit (ce qui revient à passer à une extension finie de $K$) OPS $\pi'=\pi$, alors l'image $\unsure{Z'|\mathcal Z'}$ de $\unsure{Z|\mathcal Z}$ dans $\Gamma$ est contenue dans le noyau de l'hom. $\varphi:\Gamma\longrightarrow\operatorname{Aut\,ext}(\pi)$. Or (sauf erreur) je sais prouver que cet hom. est injectif (on est ramené aussitôt au cas où $K$ est de degré de transc. $1$, et on est ramené au cas \inserted{du $\pi_1$} d'une courbe algébrique\ldots)
\margin{20}

\scanpage{22}
\uline{Théorème 2} Le foncteur $\unsure{K-\Pi_K/\Pi_{\mathbb Q}|K\to\Pi_K/\Pi_{\mathbb Q}}$ des catégories des \unsure{t.f.|v.f.} de $\mathbb Q$ vers les groupoïdes profinis sur $\Pi_{\mathbb Q}$ est \uline{fidèle}.

\margin{En fait, ce théorème n'est pas spécial à $\mathbb Q$ --- il marcherait sur un corps de base quelconque --- et en fait la dém. montre bien que c'est en fait un énoncé de nature géométrique non arithmétique (bien que pr la \unsure{démon.|démonstr.} \unsure{faut|faite|facile} une \unsure{voie} arithmétique pure)}

Si deux hom. $K\overset{f,g}{\rightrightarrows}K'$ définissent des hom de groupoïdes \inserted{sur $\Pi_{\mathbb Q}$} isomorphes
\diagram{\begin{tikzcd}[ampersand replacement=\&]
\Pi_{K'} \arrow[rr,shift left,"f^*"] \arrow[rr,shift right,"g^*"'] \arrow[dr,no head,"p'"'] \& \& \Pi_K \arrow[dl,"p"]\\
\& \Pi_{\mathbb Q}, \&
\end{tikzcd}}
(i.e. $\exists$ isomorphisme de foncteurs $\alpha:f^*\simeq g^*$, tel que $\forall\bar\eta'\in\operatorname{Ob}\Pi_{K'}$, l'e\deleted{isom}\inserted{carré}
\deleted{correspondant}
\diagram{\begin{tikzcd}[ampersand replacement=\&]
pf^*(\bar\eta') \arrow[r,phantom,"{\overset{p(\alpha)}{\simeq}}" description] \arrow[d,"\wr"'] \& pg^*(\eta) \arrow[d,no head,"\wr"]\\
p'(\bar\eta') \arrow[r,phantom,"\simeq" description] \& p'(\bar\eta')
\end{tikzcd}}
est commutatif) alors $f=g$.

L'hypothèse sur $f,g$ signifie aussi, en termes d'une clôture alg. choisie $\bar K'$ \deleted{\unsure{qui ne|qu'on} \unsure{peut}}\inserted{donnant via $\unsure{f|f'},g$ deux clôtures alg. de $\unsure{\bar K_\varepsilon|\bar K_f|\bar K}$, $\unsure{K|\widetilde K}$ de $K$} de $K'$, que \deleted{les deux hom d'exten}\inserted{l'on peut trouver un isomorphisme entre celles-ci (permettant donc d'identifier $\mathcal E_{\bar K/K}$ et $\mathcal E_{\widetilde K/K}$),} \deleted{bien que\unsure{,|.}}\inserted{de telle façon que les deux hom}
\margin{induisant ``l'identité'' sur les clôtures alg. \unsure{\deleted{$\overline{\mathbb Q}$}|$\overline{\mathbb Q}$} et $\widetilde{\mathbb Q}$ de $\mathbb Q$ dans $\bar K$ et $\widetilde K$, identifiées à celle dans $\bar K'$}
\[
f^*,g^*:\mathcal E_{\bar K'/K'}\substack{\longrightarrow\\\unsure{-|\longrightarrow}}\mathcal E_{\bar K/K}
\]
sont \unsure{égaux|egaux}, \deleted{Il \unsure{s'agit|faut} en \unsure{conclure|conclusion}}
% The following cancelled replacement is written above the preceding cancelled run.
\inserted{\deleted{\unsure{C'est} \unsure{sans} doute plus}}
\unsure{Ou bien} en termes d'une clôture alg. $\overline{\mathbb Q}$ fixée, de $\mathbb Q$, en disant que les deux hom $f^*,g^*:\pi_{K'}\rightrightarrows\pi_K$ de groupes profinis extérieurs (avec opérateurs $\Gamma_{\overline{\mathbb Q}/\mathbb Q}$) sont égaux.
\margin{On doit se ramener au cas $k=k'$ --- c'est immédiat en \deleted{\unsure{faisant le chgt de base|faire d'extension}}\inserted{\unsure{pr une|prenant|par une}} extension finie de $\mathbb Q$}
\margin{21}

\scanpage{23}
\mspage{11}
Écrivons comme tantôt $K=\varinjlim A_i$ donc \deleted{$K\otimes$} $\eta=\operatorname{Spec}(K)=\varprojlim U_i$, on a (en termes d'un pt géom. quelconque $\bar\eta$ de $\operatorname{Spec}K$ i.e. en termes d'un $\bar K$)
\[
\pi_K=\varprojlim_i\pi_1(\bar U_i{}_{\deleted{\unsure{\bar\eta|\bar K|\bar Q}}},\unsure{\bar\eta|\bar K}),\quad\text{où }\bar U_i=U_i\otimes_{\unsure{k|K}}\overline{\mathbb Q}
\]
et il suffit de voir que $\forall i,\ \overset{f_i}{\overline{f|A_i}}=\overset{g_i}{\overline{g|A_i}}$ \unsure{:|.} moyennant le fait que
\[
\pi_1(f_i^*)=\pi_1(g_i^*):\pi_{K'}\unsure{\longrightarrow|-}\pi_1(U_i)\deleted{\unsure{K|\bar\eta}}
\]
(comme hom de groupes extérieurs. Or pour $i$ fixé, on a $K'=\unsure{\varinjlim|\underset{-}{\lim}} A'_j$, où les $A'_j$ contiennent $f_i(A_i)$ et $g_i(A_i)$, donc
\[
\pi_{K'}=\varprojlim_j\pi_1(\bar V_j,\bar\eta'),\quad\text{avec }\bar V_j=\operatorname{Spec}(A_j)\otimes_k\overline{\mathbb Q}.
\]
Notons (prenant les $V_j$ \deleted{\unsure{normaux}}\inserted{réguliers}) que les hom de transition dans le syst. projectif des $\pi_1(\bar V_j,\bar\eta')$ sont surjectifs --- donc \unsure{:|..} $\pi_{K'}\longrightarrow\pi_1(\bar V_j,\bar\eta')$ est surjectif, ce qui implique que l'égalité de $f^*$ et $g^*:\pi_{K'}\substack{\unsure{-|\longrightarrow}\\\longrightarrow}\pi_{\deleted{\unsure{K|k}}}(\bar U_i)$ (comme hom extérieurs) implique celle des $\pi_1(\bar V_j)\longrightarrow\pi_1(\bar U_i)$.
Donc l'égalité $f_i=g_i$ (d'où $f=g$) sera consé-quence du résultat plus général\unsure{,|;}
\margin{22}

\scanpage{24}
\uline{\unsure{Corollaire|Corollaires}} ``purement \unsure{géométrique|géométriques}'':

\uline{Corollaire} 1 Soient $X,Y$ des schémas \inserted{de t.f. réduits connexes} sur un corps alg. clos, $k$, et $f,g:X\longrightarrow Y$ deux morphismes, on suppose que $\pi_1(f),\pi_1(g):\pi_1(X)\rightrightarrows\pi_1(Y)$ sont égaux (en tant qu'hom. extérieurs). Alors

a) Si $Y$ \deleted{est \unsure{une|un} extension \unsure{d'une|d'un}} \inserted{se plonge par un $i:Y\unsure{-|\longrightarrow}G$ un groupe algébrique commutatif $G$} extension d'une VA par un tore, il existe un $a\in Y$ (unique) tel que $g(x)=f(x)+a$ $\forall x\in X(k)$, i.e. $(i\circ g)=\tau_a\circ(i\circ f)$ ($\tau_a$ la translation par $a$).

b) Si $Y$ est une variété élémentaire d'\unsure{Artin} ou \deleted{\unsure{une}} \unsure{fibrée|fibrées} \unsure{successive|successives} des \unsure{courbes} anabéliennes, \inserted{et $X$ \deleted{normal} \inserted{\unsure{irred|(irred)}} et si $f$ ou $g$ est dominant} alors $f=g$.

\margin{\uuline{b)} \unsure{\deleted{Ce cas n'est prouvé qu'en car. $0$ \unsure{pour les courbes}; en car. \unsure{$p>0$}, on ne sait \unsure{rien sur la} structure du $\pi_1$ (\unsure{pro-$p$}, \unsure{pro-prime à $p$}), \unsure{et même} les centres des $\pi_1$\ldots}}}
\margin{Une conjecture forte établirait l'unicité des morphismes, avec $X$ irréductible et \unsure{quels} \illegible{1} applications $f,g$, \unsure{sans les supposer|faut les supposer} \unsure{dominantes|dominant} \unsure{$k$|$K$}.}
\margin{En fait dans a) il suffit de supposer que $H_1(f,\mathbb Z_\ell)=H_1(g,\mathbb Z_\ell)$ pour un $\ell$ premier à la car. \deleted{\unsure{exp.}}}

\uline{Démonstration} a) L'unicité de $a$ est claire -- et \deleted{il} \inserted{l'existence} revient, \inserted{remplaçant $Y$ par $G$ et} en choisissant une ``origine'' $a\in X(k)$, et en se ramenant au cas $f(a)=g(a)=0$, à prouver qu'alors $f=g$. D'ailleurs l'hyp. $\pi_1(f)=\pi_1(g)$ peut être affaiblie en l'hyp. que \deleted{\unsure{$f_1$|$\pi_1$|$f$} est} $H_1(f,\mathbb Z_\ell)=H_1(g,\mathbb Z_\ell)$, \inserted{$\ell$ premier à car.} -- \unsure{i.e.|ici} il suffit
\margin{23}

\scanpage{25}
\mspage{12}
d'examiner les actions de $\pi_1(f),\pi(g)$ sur les groupes abélianisés des $\pi_1$, et même sur leurs composantes $\ell$-adiques. Prenant alors jacobienne généralisée du type ``ext. d'une V.A. par un tore'' de $X$, on sait que

1°) Les morphismes $f:X\longrightarrow G$ tels que $f(a)=0$ se factorisent de façon unique par
\diagram{\begin{tikzcd}[ampersand replacement=\&]X \arrow[r,"{\mathrm{can}}"] \& J \arrow[r,"\varphi"] \& G\end{tikzcd}},
avec $\varphi$ un hom de groupes algébriques

2°) \deleted{$\varphi$ est} Un tel hom $\varphi$ est connu quand on connaît son action sur les $H_1(\ ,\mathbb Z_\ell)$ Ce qui revient à le connaître sur les points d'ordre $\ell^\nu$, quel que soit $\nu$ -- or ceux-ci sont denses\ldots

3°) $H_{\unsure{1|n}}(X,\mathbb Z_\ell)\unsure{\simeq|\overset{\sim}{\longrightarrow}}H_1(J,\mathbb Z_\ell)$

De ceci, \deleted{le \unsure{cas}} on conclut \inserted{(par 3°)} que $H_1(f)=H_1(g)$ implique (si $f=\varphi\circ\mathrm{can}$, $g=\psi\circ\mathrm{can}$) $H_{\unsure{1|n}}(\varphi)=H_1(\psi)$, donc par 2°) que $\varphi=\psi$, donc $f=g$, cqfd

\deleted{b)} Notons que l'on ne peut conclure de l'hyp. \inserted{de a)} que $f=g$, puisque pour $X=Y=G$, les translations $\tau_a$ induisent l'identité sur $\pi_1(G)$ (du moins \unsure{en|\deleted{en}|\deleted{si}} la car. $0$ ou si $G$ est une VA), sans être l'identité.

b) On va pourtant prouver l'égalité,
\margin{24}

\scanpage{26}
sous l'hypothèse anabélienne faite, et celle de dominance. (Cette deuxième n'est pas superflue, comme on voit en prenant \deleted{\unsure{pour} $Y$ \inserted{$=G\backslash Z$, où $G$ est} un groupe elliptique et $Z$ un sous-groupe fini $\unsure{+1|\ne 1}$, et pour $f,g$} \inserted{deux applications constantes de $X$ dans $Y$!})

Tout d'abord, on vérifie qu'une variété élémentaire d'Artin se plonge bien \inserted{par un $i:Y\longrightarrow G$} dans une ext. d'une V.A. par un tore -- donc on aura $(ig)=\tau_u(i\circ f)$. \deleted{Donc \unsure{tout|faut}} \inserted{D'autre part,} \deleted{Quitte à remplacer} $f$,\deleted{$g$} étant dominante, il existe une partie ouverte \inserted{non vide} $\unsure{V|W}\subset Y$ contenue dans $\operatorname{Im}f$, et on aura donc $\unsure{V|W}+u\subset Y$. Soit $U\subset X$ un ouvert non vide tel que $f(U)\subset V$ \deleted{$g(U)\subset V$} \deleted{\unsure{donc qui} $\unsure{Q}$} quitte à remplacer $X$ par $U$, on est ramené au cas où $f$\deleted{,$g$} envoie $X$ dans un ouvert $V$ de $Y$ tel que $V+u\subset Y$, \unsure{on|et|ou} $g\circ i=\tau_u(i\circ f)$. On \deleted{\unsure{sait}} \inserted{voit} que $\pi_1(f)$ est un composé (où $f':X\unsure{-|\longrightarrow}U$ est \unsure{\deleted{def.}|def.} induit par $f$ et $j:U\unsure{-|\longrightarrow}Y$ est l'inclusion)
% Each empty editorial doubt marker below a dashed shaft concerns only its terminal head: dash first, right arrow alternative. Objects and source labels are not thereby made doubtful.
\diagram{\begin{tikzcd}[ampersand replacement=\&]\pi_1(X) \arrow[r,dash,"\pi_1(f')","{\unsure{}}"'] \& \pi_1(U) \arrow[r,"\pi_1(j)"] \& \pi_1(Y)\end{tikzcd}}
tandis que $\pi_1(g)$ est le composé
\diagram{\begin{tikzcd}[ampersand replacement=\&]\pi_1(X) \arrow[r,dash,"\pi_1(f')","{\unsure{}}"'] \& \pi_1(U) \arrow[r,dash,"{\pi_1(\unsure{j'|g'})}","{\unsure{}}"'] \& \pi_1(Y)\end{tikzcd}}
% The overwritten j'/g' label is one unresolved glyph; the visible overtracing does not securely distinguish a cancelled g from a replacement j.
où $\deleted{\unsure{j|g}}$ $j':U\unsure{-|\longrightarrow}Y$ est induit par $\tau_u$.
\margin{25}

\scanpage{27}
\mspage{13}
L'hyp. que $\pi_1(f')$ \deleted{et} $=\pi_1(g)$ signifie donc qu'il existe $\alpha\in\pi_1(Y)$, tel que
\[
\pi_1(j')(\unsure{\gamma|\sigma})=\inserted{\operatorname{int}(\alpha)}\pi_1(j)(\unsure{\gamma|\sigma})
\]
pour tt $\unsure{\gamma|\sigma}\in\operatorname{Im}(\pi_1(X)\overset{\pi_1(f)}{\unsure{-|\longrightarrow}}\pi_1(U))$. Or cette image est un sous-groupe \uline{ouvert} de $\pi_1(U)$ (tjrs vrai pour un morphisme dominant!). \deleted{D'ailleurs} Donc on est ramené à ceci: Soit $U$ ouvert \inserted{$\ne\varnothing$} de $Y$, et $u\in G$, tels que $\tau_uU\subset Y$ [et tels que (désignant par $j,j'$ les \deleted{applications} \inserted{morphismes} $y\longrightarrow y$ et $y\unsure{\mapstochar\relbar|\longmapsto}y+u$ de $U$ dans $Y$) $\pi_1(\unsure{f|j})$ et $\pi_1(\unsure{f'|j'})$ coïncident extérieurement sur un sous-groupe ouvert de $\pi_1(U)$] -- \uline{alors} $j=j'$ \unsure{i.e|v.c} $u=0$.

\margin{\uline{NB} Ceci n'est pas plus fort que le corollaire, comme on voit en prenant pour $X$ le rev. étale fini de $U$ correspondant \unsure{au|à un} sous-groupe ouvert envisagé de $\pi_1(U)$\unsure{. --}}

Finalement, je m'aperçois que je n'arrive pas à le prouver par voie géométrique -- pour ceci je me rabats sur une méthode \uline{arithmétique};
\margin{Si -- cf. plus bas!}

\deleted{\unsure{Je pense|Je ferai} d'ailleurs \deleted{\unsure{affaibli}} \unsure{généraliser|généralisé} le th. géométrique; on montre que (quitte à remplacer $X$ par un ouvert non vide \unsure{convenable|connexe} ce qui ne détruit pas les hyp ni ne modifie la conclusion) OPS que $X$ est une variété}
\margin{26}
